\section{Conclusiones}

El desarrollo del brazo robótico fue una experiencia muy completa para el equipo, ya que permitió poner en práctica muchos de los conocimientos que fuimos adquiriendo a lo largo de la carrera. El resultado final fue un prototipo completamente funcional, que puede moverse mediante cuatro grados de libertad y que permite ser controlado tanto desde el comando espejo como desde la interfaz gráfica desarrollada. Durante las pruebas, el sistema demostró un buen funcionamiento al manipular objetos metálicos livianos como tornillos, tuercas y pequeños componentes, cumpliendo con los objetivos principales del proyecto.

Uno de los aspectos más satisfactorios del proyecto fue poder ver el brazo funcionando correctamente luego de todo el proceso de diseño, prueba y corrección. En especial, el correcto funcionamiento de la grabación y reproducción de movimientos, que al principio era una de las partes que más dudas generaba, terminó funcionando incluso mejor de lo esperado, lo que fue una gran motivación para todo el equipo.

La parte más compleja del proyecto fue el diseño de la estructura mecánica. A diferencia de la electrónica o la programación, donde existe mucha información y ejemplos disponibles, la estructura requirió un proceso mucho más experimental. Se buscó lograr una estructura que fuera sencilla de armar y desarmar, pensada especialmente para ser utilizada en entornos educativos, y que al mismo tiempo fuera resistente y visualmente atractiva. Para esto se diseñaron encastres impresos en 3D que permitieran reducir el uso de tornillos y facilitar el armado, lo que además hizo que el montaje del brazo fuera una experiencia entretenida y didáctica.

Aunque el prototipo funciona correctamente, se identificaron varias limitaciones técnicas. Los servomotores utilizados presentan una pequeña zona muerta, que si bien es pequeña, afecta la precisión del posicionamiento, especialmente cuando el brazo se encuentra extendido. Además, al no contar con sensores de posición absoluta, fue necesario implementar una rutina de inicio que lleva a todos los servos a una posición segura, lo que genera movimientos bruscos y esfuerzos sobre la estructura. Estas limitaciones no impiden el uso del sistema, pero muestran claramente aspectos que podrían mejorarse en futuras versiones.

Desde el punto de vista del aprendizaje, el proyecto fue muy importante para el equipo. Todos los integrantes participaron en las distintas etapas, por lo que se fortalecieron conocimientos en diseño mecánico, electrónica, programación y, especialmente, en la integración de sistemas. Se pudo comprobar que los contenidos vistos en la carrera no quedan solo en la teoría, sino que se pueden aplicar directamente en proyectos reales y funcionales.

El trabajo en equipo fue fundamental para poder completar el proyecto. Si bien aparecieron diferencias durante el proceso, se lograron resolver de forma adecuada, manteniendo siempre el objetivo común. Esto permitió que el grupo funcionara de forma organizada y que se pudiera avanzar de manera constante, logrando un muy buen resultado final.

En cuanto a su uso educativo, se considera que el brazo robótico es totalmente viable como herramienta didáctica. No solo permite mostrar el funcionamiento de los sistemas robóticos, sino que también ofrece la posibilidad de que otros estudiantes lo modifiquen, lo mejoren o le agreguen nuevas funciones, como una garra mecánica o un electroimán más potente. Esto lo convierte en una base muy interesante para futuros trabajos prácticos.

Por último, el proyecto deja abiertas varias posibilidades de mejora. Entre ellas se encuentra la creación de una versión instalable de la interfaz gráfica para facilitar su uso por personas con menos experiencia técnica, el diseño de una placa de circuito impreso que reduzca el cableado y mejore el orden general del sistema, y la implementación de un efector de garra que quedó pendiente por falta de tiempo. Más allá de estas mejoras, el proyecto logró su objetivo principal: desarrollar un brazo robótico didáctico, funcional, replicable y con un fuerte valor educativo, dejando una base sólida para futuros proyectos.

% TODO: Aspectos positivos
% TODO: Aspectos negativoss
% TODO: Si fueramos a distribuirlo como kit... el circuito deberia incorporarse en una PCB o shield para el arduino, el tablero debería ser incorporado de manera similar o ser una unidad prefabricada.

% TODO: Mejoras a futuro